% Case C: Language Switching

\footnotesize

\begin{frame}{Case C：Language Switching — 设计意图与总览}
\scriptsize
\textbf{核心假设：}NLA 能从 layer-20 activation 中区分当前输出语言、语言切换计划，以及文化 cue 触发的语言语义邻域。

\textbf{实验设计：}覆盖 English-only、Chinese-only、显式 SWITCH、English-only 但含中国地名、英文问题要求中文回答；另用 R1 code-switch 作为早期对照。

\textbf{来源：}R3 \code{language\_switching} + R1 Plan D 多语言输入。

\vspace{0.25em}
\begin{tabularx}{\linewidth}{lrrrrX}
\toprule
条件 & hit & cos & FVU & delta & 解读 \\
\midrule
C1 English-only & 3.05\% & 0.944 & 0.223 & 0.356 & 低 false positive baseline \\
C2 Chinese-only & 49.62\% & 0.952 & 0.190 & 0.478 & 中文任务，最高 cos \\
C3 SWITCH & 44.35\% & 0.905 & 0.321 & 0.341 & 切换任务强 language hit \\
C4 implicit language & 61.83\% & 0.941 & 0.233 & 0.268 & 地名/文化 cue 触发最高 hit \\
C5 English Q Chinese A & 52.05\% & 0.890 & 0.449 & 0.762 & 语言错配，强中文信号但低保真 \\
\bottomrule
\end{tabularx}

\vspace{0.3em}
\alert{总览：}语言是本轮最稳定的可读方向；但 C4 说明 language hit 不一定等于"即将切换语言"。
\end{frame}

\begin{frame}{Case C：C3 SWITCH（显式切换例）}
\scriptsize
\textbf{来源：}R3 \code{C3\_switch\_after\_marker}。

\textbf{输入 Prompt}
\texttt{"Explain why tea is popular in many cultures in English first. After the exact word SWITCH, continue the same explanation in Chinese."}

\textbf{Response:}

 \dots cultures due to its unique flavor,\colorbox{green!15}{ health}\textsuperscript{[{13}]} benefits, and social significance. It \dots making it\colorbox{green!15}{ an}\textsuperscript{[{80}]} integral part of daily life in many societies.

SW\colorbox{green!15}{ITCH}\textsuperscript{[{91}]}

茶在许多文化中流行 \dots

\textbf{结果：}\code{cos = 0.905}, \code{FVU = 0.321}，language/translation hit=\textbf{44.35\%}。

\textbf{NLA decode：}见下一页，按 \code{pos 13/80/91} 展示 \code{scores.csv} 原始 explanation 全文。

\alert{对照解释：}C3 证明语言切换相关 activation 可读；但它不能单独证明提前检测，因为命中大量发生在 SWITCH 和中文输出附近。
\end{frame}

\begin{frame}{Case C：C3 SWITCH — NLA decode 原文}
\tiny
\textbf{来源：}R3 \code{C3\_switch\_after\_marker}，对应上一页高亮 token。

\begin{tabularx}{\linewidth}{l>{\raggedright\arraybackslash}X>{\raggedright\arraybackslash}p{2.0cm}}
\toprule
位置 & NLA decode 原始文本 & 备注 \\
\midrule
pos 13 & Structured article format with formal descriptive tone introducing ``Green tea,'' establishing cultural context and popularity in Japan. & English 区段：茶文化语义 \\
pos 80 & Structured informational format with formal English language article tone, detailing tea's cultural significance and health benefits, establishing a pattern of explaining social habits. & SWITCH 前后，仍是茶文化主题 \\
pos 91 & Structured Chinese language article format with numbered sections and translations, presenting a video's content about ``Happy Birthday'' in English and Chinese with formal definitions. & \textbf{明确双语/中文命中} \\
\bottomrule
\end{tabularx}
\end{frame}

\begin{frame}{Case C：C1/C2 紧凑对照}
\scriptsize
\textbf{来源：}R3 \code{C1\_english\_only}, \code{C2\_chinese\_only}。

\begin{columns}[T,onlytextwidth]
\begin{column}{0.48\linewidth}
\textbf{C1 English-only}

\textbf{Prompt:} \texttt{"Explain why tea is popular in many cultures. Answer only in English."}

\textbf{Response:}

Tea is popular in many cultures due\colorbox{green!15}{ to}\textsuperscript{[{8}]} its unique flavor \dots

\vspace{1em}
\textbf{NLA decode：}pos 8: Structured informational article format with formal tone establishes educational content about tea's cultural significance and popularity；language hit=\textbf{3.05\%}。

\vspace{1em}
\textbf{结果：}\code{cos = 0.944}。
\end{column}
\vrule{}\hspace{0.01\linewidth}
\begin{column}{0.48\linewidth}
\textbf{C2 Chinese-only}

\textbf{Prompt:} \texttt{"请用中文解释为什么茶在许多文化中很受欢迎。"}

\textbf{Response:}

茶在许多文化中非常\colorbox{green!15}{受欢迎}\textsuperscript{[{6}]} \dots

\vspace{1em}
\textbf{NLA decode：}pos 6: Structured article format with formal introduction establishing ``茶''的多方面魅力，文章采用中文标题和英文解析方式，预设内容将详细阐述茶的受欢迎原因。

\vspace{1em}
\textbf{结果：}\code{cos = 0.952}, language hit=\textbf{49.62\%}。
\end{column}
\end{columns}

\vspace{0.3em}
\alert{对照解释：}C1 给出 English-only 的低命中基线；C2 显示同一主题换成中文后，语言表征强且高保真。
\end{frame}

\begin{frame}{Case C：C4/C5 紧凑对照}
\scriptsize
\textbf{来源：}R3 \code{C4\_implicit\_language}, \code{C5\_malformed\_pair\_sim}。

\begin{columns}[T,onlytextwidth]
\begin{column}{0.48\linewidth}
\textbf{C4 Implicit Language}

\textbf{Prompt:} \texttt{"Explain the cultural importance of tea in Hangzhou and Taipei. Use English only."}

\textbf{Response:}

Tea holds significant cultural importance in both Hang\colorbox{green!15}{zhou}\textsuperscript{[{9}]} \dots

\vspace{1em}
\textbf{NLA decode：}pos 9: Structured article format with formal descriptive tone establishes an informational article about Chinese tea culture, introducing a cultural landmark with historical context；最高 language hit=\textbf{61.83\%}。

\vspace{1em}
\textbf{结果：}\code{cos = 0.941}。
\end{column}
\vrule{}\hspace{0.01\linewidth}
\begin{column}{0.48\linewidth}
\textbf{C5 English Q, Chinese A}

\textbf{Prompt:} \texttt{"Answer the following English question in Chinese ... Why is tea popular in many cultures?"}

\textbf{Response:}

茶在许多文化中\colorbox{green!15}{流行}\textsuperscript{[{5}]}，是因为 \dots

\vspace{1em}
\textbf{NLA decode：}pos 5: Structured format with numbered definitions and translations establishes a pattern of presenting ``Popular'' as a common English word with its Chinese translation and explanation。

\vspace{1em}
\textbf{结果：}\code{cos = 0.890}, language hit=\textbf{52.05\%}。
\end{column}
\end{columns}

\vspace{0.4em}
\alert{关键反证：}C4 要求 English-only，却因为 Hangzhou/Taipei 产生最高 language hit。后续若要证明"切换前"信号，必须做 lead-time：只统计第一个中文 token 之前的 NLA 命中。
\end{frame}

\begin{frame}{Case C：R1 Code-Switch 早期对照}
\scriptsize
\textbf{来源：}R1 Plan D / \code{D3\_mixed\_code\_switch}。

\textbf{输入 Prompt}

\textbf{System:} \texttt{"You are Qwen, created by Alibaba Cloud. You are a helpful assistant."}

\textbf{User:} \texttt{"Explain AI first in Chinese, then in English"}

\textbf{Response:}

当然\colorbox{green!15}{可以}\textsuperscript{[{1}]}。

首先\colorbox{green!15}{用}\textsuperscript{[{4}]}中文解释AI：人工智能 \dots

\begin{tabular}{lp{6.2cm}l}
\toprule
位置 & NLA decode & 备注 \\
\midrule
pos 0--3 & \texttt{"当然"}, \texttt{"首先用中文"}, \texttt{"那么用中文来解释一下"}, \texttt{"首先，我先用中文简要解释一下"} & 读到语言规划过程 \\
pos 4--15 & \texttt{"首先，让我们先用中文表达这个概念"}；\texttt{"我们常说的'人工智能'是什么意思？以下简单介绍一下："} & 持续追踪"先中后英"组织方式 \\
\bottomrule
\end{tabular}

\textbf{结果：}\code{cos = 0.886}, \code{MSE = 0.229}, \code{norm = 118.3}。

\alert{对照解释：}R1 D3 的 decode 最像"语言规划"本身，而 R3 C4 提醒我们不能把所有 language hit 都解释成 switch planning。两者合起来给出更谨慎的结论：语言表征强可读，提前切换证据仍需单独验证。
\end{frame}
